\subsection{2026-09-21 --- Feature iteration, first pass: 0.707 $\to$ 0.808}
\label{sec:feature-iter-1}

Seven cumulative stages, each scored with the full rolling-origin protocol
(Section~\ref{sec:cv-harness}). Notebook: \texttt{src/modeling/feature\_iteration.ipynb};
builders in \texttt{src/features.py}. Total runtime 30 minutes (28 training runs).

\begin{table}[htbp]
  \centering
  \caption{Cumulative feature stages. $\Delta$ is versus the previous row; ``folds up'' counts how
  many of the four validation years improved.}
  \label{tab:feature-iter-1}
  \resizebox{\linewidth}{!}{%
  \begin{tabular}{llrrrrrrrr}
    \toprule
    Stage & & \#feat & 2019 & 2020 & 2021 & 2022 & Mean & $\Delta$ & Folds up \\
    \midrule
    00 & baseline (raw)      &  54 & 0.799 & 0.679 & 0.666 & 0.686 & 0.707 & --- & --- \\
    01 & seasonal percentile &  58 & 0.782 & 0.673 & 0.643 & 0.651 & 0.687 & $-$0.020 & \textbf{0/4} \\
    09 & inflow regime       &  62 & 0.783 & 0.668 & 0.679 & 0.641 & 0.693 & $+$0.006 & 2/4 \\
    02 & price, case-relative&  70 & 0.751 & 0.687 & 0.671 & 0.686 & 0.699 & $+$0.006 & 2/4 \\
    06 & \textbf{price $-$ water value spread} & 81 & 0.790 & 0.721 & 0.782 & 0.802 & 0.774 & $\mathbf{+0.075}$ & \textbf{4/4} \\
    07 & \textbf{terminal water value} & 85 & 0.831 & 0.745 & 0.808 & 0.845 & 0.807 & $\mathbf{+0.034}$ & \textbf{4/4} \\
    10--12 & physics (head, start cost) & 157 & 0.828 & 0.761 & 0.810 & 0.834 & 0.808 & $+$0.001 & 2/4 \\
    08 & drop dead columns   & 134 & 0.824 & 0.772 & 0.813 & 0.823 & 0.808 & $-$0.000 & 2/4 \\
    \bottomrule
  \end{tabular}}
\end{table}

\finding{\textbf{Overall $+0.101$, from 0.707 to 0.808}, and essentially all of it comes from two
stages: the price\,$-$\,water-value spread ($+0.075$) and the terminal water value ($+0.034$).
Both improved \textbf{all four folds}, which is the criterion we set in
Section~\ref{sec:cv-harness}. Fold SD also fell from 0.061 to 0.025, so the model became more
stable across years as well as better.}

\finding{\textbf{Univariate correlation was a poor predictor of feature value --- in one case
actively misleading.} The seasonal storage percentile had the strongest correlation measured
anywhere in the project ($+0.537$, Section~\ref{sec:eda-round-3b}) and it made the model
\emph{worse}: $-0.020$ mean, \textbf{0 of 4 folds improved}, and it accounts for 0.9\% of total
model gain. The case-relative price z-score ($+0.263$ correlation) added $+0.006$ on 2 of 4 folds.
Meanwhile the raw spread, whose correlation was a middling $+0.241$, contributed
\textbf{60.8\% of all model gain}. The caveat recorded in Section~\ref{sec:cv-harness} --- that
correlation ignores what the other features already supply --- turned out to be the dominant
effect, not a footnote.}

\finding{\textbf{The metric rewards \emph{within-case} discrimination, which explains the
pattern.} Micro-average AUC pools every (case, hour, unit) prediction, so most of the ranking
problem is deciding \emph{which hours within a week} are ON. A case-level scalar can only shift
whole cases up or down. Measured on the fitted model: \textbf{hourly-varying features hold 96.4\%
of total gain; case-level constants hold 3.6\%}. The seasonal percentile, inflow regime and
storage anomaly are all case-level constants --- which is why they underperformed their
correlations. Prioritise features that vary within the case.}

\finding{\textbf{The single most important feature is an economic decision rule, not a
statistical transform.} \texttt{f06\_spread\_sys\_positive} --- the binary ``is price above the
system water value'' --- has a gain of 207{,}444, roughly \textbf{14$\times$ the next feature}. That
is literally the optimizer's dispatch condition. Start-cost economics
(\texttt{f10\_breakeven\_hours\_*}) place five entries in the top 25. This supports building
features that mirror what SHOP computes, rather than more transformations of the inputs.}

\finding{\textbf{Gain share is not incremental value.} The physics stage contributed 72 columns
and 13.5\% of model gain, but only $+0.0009$ mean AUC on 2 of 4 folds --- the model uses those
columns, yet they add almost nothing beyond what the spread features already carry. Stage-level
deltas and importance rankings answer different questions; neither alone is sufficient.}

\subsubsection*{Bookkeeping.} Dropping the 23 single-valued columns was correctly neutral
($-0.0002$), confirming the model had already ignored them; it did cut fold SD from 0.033 to
0.025. Twenty-one of those 23 were generated by the physics stage itself --- static penstock
counts, and Haukeli's entire physics block, which is constant because its reservoir (Vatjern) is
one of the two with no historical record (Section~\ref{sec:eda-round-2}, finding 6).

\subsubsection*{Next.} A screening mode has been added to the harness
(\texttt{screening\_eval}): one training run on 2015--2019, scored on 2020/2021/2022, 150 trees,
every second training hour. Validation is inference and therefore nearly free, so three
validation years cost the same as one --- the four-fold expense came from four \emph{trainings}.
Expected $\sim$16$\times$ faster, making the remaining $\sim$13 themes practical to test in bulk.
First candidates: the water-budget family (available energy vs.\ profitable hours, and the
greedy merit-dispatch margin), which extends the winning economic-rule pattern, and the
global-vs-expanding A/B on the seasonal percentiles.
